\subsection{理想}

定义 4. 设 A 是一个环。A 的一个子集 \(\mathfrak{a}\) 称为 A 的左（相应地，右）理想，如果它是 A 的加法群的子群，并且关系 \(a \in \mathbf{A}\) ， \(x \in \mathfrak{a}\) 蕴含 \(ax \in \mathfrak{a}\) （相应地， \(xa \in \mathfrak{a}\)）。\(\mathfrak{a}\) 称为 A 的双侧理想，如果它既是 A 的左理想又是 A 的右理想。

左理想的定义可以用以下关系表示

\[
0 \in a, \quad a + a \subset a, \quad A. a \subset a
\]

关系 \(-a \subset a\) 由公式 \((-1).x = -x\) 和 \(A.a \subset a\) 得出。对所有 \(x \in A\) ，设 \(\gamma_x\) 为映射 \(a \mapsto xa\) （从 \(A\) 到 \(A\) 中）；作用 \(x \mapsto \gamma_x\) 给加法群 \(A^+\) （\(A\) 的）赋予一个以 \(A\) 为算子集的带算子群的结构。A 的左理想正是 \(A\) 的加法群 \(A^+\) 中在该作用下稳定的子群。

环 A 中的左理想恰好就是相反环 \(A^{0}\) 的右理想。在交换环中，这三种理想是相同的；它们简称为理想。

例子。(1) 设 A 是一个环。集合 A 是 A 的一个双边理想；由 0 组成的集合也是，它被称为零理想，有时记为 0 或 (0)，而非 \{0\}。

\begin{enumerate}
\def\labelenumi{(\arabic{enumi})}
\setcounter{enumi}{1}
\item
  对每个元素 \(a\) （\(\mathbf{A}\) 的），集合 \(\mathbf{A}.a\) （即 \(a\) 的左倍数组成的集合）是一个左理想；类似地，集合 \(a.\mathbf{A}\) 是一个右理想。当 \(a\) 属于 \(\mathbf{A}\) 的中心时， \(\mathbf{A}.a = a.\mathbf{A}\) ；这个理想称为由 \(a\) 生成的主理想，记作 (a)。(a) = A 当且仅当 \(a\) 是可逆的。
\item
  设 M 是 A 的一个子集。由元素 \(x \in A\) （满足对所有 \(y \in M\) 有 xy = 0）组成的集合是 A 的一个左理想，称为 M 的左零化子。M 的右零化子类似地定义。
\item
  A 的左（相应地，右，双侧）理想的交仍是左（相应地，右，双侧）理想。因此，给定 A 的子集 X，存在包含 X 的最小左（相应地，右，双侧）理想；它被称为由 X 生成的左（相应地，右，双侧）理想。
\end{enumerate}

设 \(\mathfrak{a}\) 为 \(\mathbf{A}\) 的一个左理想。条件 \(1 \notin \mathfrak{a}, \mathfrak{a} \neq \mathbf{A}\) 显然是等价的。

定义 5. 设 A 是一个环。以语言上的滥用，称一个左理想 a 为极大的，如果它是不同于 A 的左理想组成的集合中的一个极大元。

换句话说，a 是极大的，如果 \(a \neq A\) 且 A 中包含 a 的仅有的左理想是 a 和 A。

定理 1（Krull）。设 A 是一个环，a 是 A 的一个不同于 A 的左理想。则存在 A 的一个包含 a 的极大理想 m。

把 A 视为以左乘法作用于 A 的加法群 \(A^{+}\) 。于是 A 的左理想就是 \(A^{+}\) 的稳定子群。因此，该定理由将 § 4 第 3 号的命题 3 应用于子集 \(P = \{1\}\) （\(A^{+}\) 的）而得出。

命题 3. 设 A 是一个环， \((x_{\lambda})_{\lambda \in L}\) 是 A 中元素的一个族，且 a（相应地 b）为和 \(\sum_{\lambda \in L} a_{\lambda} x_{\lambda}\) （其中 \((a_{\lambda})_{\lambda \in L}\) 是 A 中元素的一个有限支撑族）（相应地，和 \(\sum_{\lambda \in L} a_{\lambda} x_{\lambda} b_{\lambda}\) ，其中 \((a_{\lambda})_{\lambda \in L}, (b_{\lambda})_{\lambda \in L}\) 是 A 中元素的有限支撑族）组成的集合。则 a（相应地 b）是由元素 \(x_{\lambda}\) 生成的 A 的左（相应地，双边）理想。

公式

\[
0 = \sum_ {\lambda \in L} 0. x _ {\lambda}\tag{18}
\]

\[
\sum_ {\lambda \in L} a _ {\lambda} x _ {\lambda} + \sum_ {\lambda \in L} a _ {\lambda} ^ {\prime} x _ {\lambda} = \sum_ {\lambda \in L} (a _ {\lambda} + a _ {\lambda} ^ {\prime}) x _ {\lambda}\tag{19}
\]

\[
a. \sum_ {\lambda \in L} a _ {\lambda} x _ {\lambda} = \sum_ {\lambda \in L} (a a _ {\lambda}) x _ {\lambda}\tag{20}
\]

证明 \(\mathfrak{a}\) 是一个左理想。设 \(\mathfrak{a}'\) 是一个左理想，使得 \(x_{\lambda} \in \mathfrak{a}'\) 对所有 \(\lambda \in \mathbf{L}\) 成立，并设 \((a_{\lambda})_{\lambda \in L}\) 是 A 中的一个有限支撑族。则 \(a_{\lambda}x_{\lambda}\in \mathfrak{a}'\) 对所有 \(\lambda \in L\) 成立，因此 \(\sum_{\lambda \in L}a_{\lambda}x_{\lambda}\in \mathfrak{a}'\) ；从而 \(\mathfrak{a}\subset \mathfrak{a}'\) 。因此 \(\mathfrak{a}\) 是由 \(x_{\lambda}\) 生成的 A 的左理想。对 \(\mathfrak{b}\) 的论证是类似的。

命题 4. 设 A 是一个环，且 \((\mathfrak{a}_{\lambda})_{\lambda \in L}\) 是 A 的一族左理想。由 \(\bigcup_{\lambda \in L} \mathfrak{a}_{\lambda}\) 生成的左理想由形如 \(\sum_{\lambda \in L} y_{\lambda}\) 的和组成，其中 \((y_{\lambda})_{\lambda \in L}\) 是一个有限支撑族，满足 \(y_{\lambda} \in \mathfrak{a}_{\lambda}\) 对所有 \(\lambda \in L\) 成立。

设 \(\mathfrak{a}\) 为由和 \(\sum_{\lambda \in \mathbf{L}} y_{\lambda}\) （满足 \(y_{\lambda} \in \mathfrak{a}_{\lambda}\) ，对所有 \(\lambda \in \mathbf{L}\) ）组成的集合。公式 \(\sum_{\lambda \in \mathbf{L}} x_{\lambda} + \sum_{\lambda \in \mathbf{L}} y_{\lambda} = \sum_{\lambda \in \mathbf{L}} (x_{\lambda} + y_{\lambda})\) 和 \(a\) . \(\sum_{\lambda \in \mathbf{L}} x_{\lambda} = \sum_{\lambda \in \mathbf{L}} ax_{\lambda}\) 表明 \(\mathfrak{a}\) 是 \(\mathbf{A}\) 的一个左理想。设 \(\lambda \in \mathbf{L}\) 且 \(x \in \mathfrak{a}_{\lambda}\) ；写 \(y_{\lambda} = x\) 及 \(y_{\mu} = 0\) （对 \(\mu \neq \lambda\) ）；则 \(x = \sum_{\lambda \in \mathbf{L}} y_{\lambda}\) ，从而 \(x \in \mathfrak{a}\) ，最终 \(\mathfrak{a}_{\lambda} \subset \mathfrak{a}\) 。若左理想 \(\mathfrak{a}'\) 包含 \(\mathfrak{a}_{\lambda}\) （对所有 \(\lambda \in \mathbf{L}\) ），则它显然包含 \(\mathfrak{a}\) ，因此 \(\mathfrak{a}\) 由 \(\bigcup_{\lambda \in \mathbf{L}} \mathfrak{a}_{\lambda}\) 生成。

理想 \(\mathfrak{a}\) （由 \(\bigcup_{\lambda \in L} \mathfrak{a}_{\lambda}\) 生成）称为左理想 \(\mathfrak{a}_{\lambda}\) 的和，记作 \(\sum_{\lambda \in L} \mathfrak{a}_{\lambda}\) （参见 II，§ 1，第 7 号）。特别地，两个左理想的和 \(\mathfrak{a}_{1} + \mathfrak{a}_{2}\) 由形如 \(a_{1} + a_{2}\) 的和组成，其中 \(a_{1} \in \mathfrak{a}_{1}\) 且 \(a_{2} \in \mathfrak{a}_{2}\) 。

